\documentclass[11pt]{article}

\usepackage{amsmath,amssymb,amsthm}
\usepackage{algorithm}
\usepackage{algorithmic}
\usepackage{fullpage}
\usepackage{hyperref}

\newtheorem{theorem}{Theorem}
\newtheorem{lemma}{Lemma}
\newtheorem{definition}{Definition}
\newtheorem{assumption}{Assumption}
\newcommand{\E}{\mathbb{E}}
\newcommand{\R}{\mathbb{R}}

\begin{document}

\section*{Algorithm Name}
Clipped Adam with Momentum (gradient-norm clipping with first and second moments). For theory, we also consider the special case with no momentum and no adaptivity (pure clipped gradient descent), which is a special case obtained by setting $\beta_1=\beta_2=0$ and omitting bias-corrections.

\section*{Mathematical Formulation}
We consider minimizing a differentiable (possibly non-convex) function $f:\R^d \to \R$. Let $\nabla f(x)$ denote its gradient at $x$.

Clipping operator (global $\ell_2$-norm clipping): for a clipping threshold $C>0$,
\[
\operatorname{clip}_C(g) \;\triangleq\; \frac{g}{\max\{1,\|g\|_2/C\}} \,.
\]
Given hyperparameters $\eta>0$ (base stepsize), $\beta_1\in[0,1)$ (first-moment momentum), $\beta_2\in[0,1)$ (second-moment momentum), $\epsilon>0$ (stabilizer), the Clipped Adam with Momentum updates are:
\begin{align*}
g_t &\leftarrow \nabla f(x_t), \\
u_t &\leftarrow \operatorname{clip}_C(g_t) = \frac{g_t}{\max\{1,\|g_t\|_2/C\}}, \\
m_t &\leftarrow \beta_1 m_{t-1} + (1-\beta_1) u_t, \\
v_t &\leftarrow \beta_2 v_{t-1} + (1-\beta_2) (u_t \odot u_t) \quad (\text{elementwise square}), \\
\hat{m}_t &\leftarrow \frac{m_t}{1-\beta_1^t}, \qquad
\hat{v}_t \leftarrow \frac{v_t}{1-\beta_2^t}, \\
x_{t+1} &\leftarrow x_t - \eta \, \frac{\hat{m}_t}{\sqrt{\hat{v}_t} + \epsilon} \quad (\text{division is elementwise}).
\end{align*}
Initialization: $x_1\in\R^d$ given, $m_0=0\in\R^d$, $v_0=0\in\R^d$.

For analysis, we will also study the special case (Clipped Gradient Descent, CGD):
\[
x_{t+1} \leftarrow x_t - \eta \, u_t \quad \text{with} \quad u_t=\operatorname{clip}_C(\nabla f(x_t)),
\]
which is obtained by setting $\beta_1=\beta_2=0$ and dropping bias corrections.

\section*{Pseudocode}
\begin{algorithm}[H]
\caption{Clipped Adam with Momentum}
\begin{algorithmic}[1]
\STATE Input: initial point $x_1\in\R^d$, stepsize $\eta>0$, clipping threshold $C>0$, momentum $\beta_1\in[0,1)$, second-moment $\beta_2\in[0,1)$, stabilizer $\epsilon>0$, max iterations $T$
\STATE Initialize $m_0\leftarrow 0$, $v_0\leftarrow 0$
\FOR{$t=1$ to $T$}
  \STATE $g_t \leftarrow \nabla f(x_t)$
  \STATE $u_t \leftarrow g_t / \max\{1,\|g_t\|_2/C\}$  \hfill (global $\ell_2$-norm clipping)
  \STATE $m_t \leftarrow \beta_1 m_{t-1} + (1-\beta_1) u_t$
  \STATE $v_t \leftarrow \beta_2 v_{t-1} + (1-\beta_2) (u_t \odot u_t)$
  \STATE $\hat{m}_t \leftarrow m_t/(1-\beta_1^t)$, $\hat{v}_t \leftarrow v_t/(1-\beta_2^t)$
  \STATE $x_{t+1} \leftarrow x_t - \eta \, \hat{m}_t / (\sqrt{\hat{v}_t} + \epsilon)$ \hfill (elementwise division)
\ENDFOR
\STATE Output: $x_{T+1}$
\end{algorithmic}
\end{algorithm}

\section*{Dry Run Example}
We apply the algorithm to $f(w)=(w-3)^2$ (one dimension, $d=1$), with $w_1=0$, $\eta=0.1$, $\beta_1=0.9$, $\beta_2=0.999$, $\epsilon=10^{-8}$, $C=2$, for $3$ iterations. In $d=1$, vectors are scalars and elementwise operations reduce to scalar ones.

At step $t$:
- Gradient: $g_t = f'(w_t) = 2(w_t-3)$
- Clipping: $u_t = g_t/\max\{1,|g_t|/2\}$

Initialization: $m_0=0$, $v_0=0$.

Iteration $t=1$:
- $w_1=0$, $f(w_1)=9$, $g_1=2(0-3)=-6$, $|g_1|=6>C$
- $u_1 = -6 / (6/2) = -2$
- $m_1 = 0.9\cdot 0 + 0.1\cdot(-2) = -0.2$
- $v_1 = 0.999\cdot 0 + 0.001\cdot 4 = 0.004$
- Bias-corrected: $\hat{m}_1 = -0.2/(1-0.9) = -2$, $\hat{v}_1 = 0.004/(1-0.999) = 4$
- Update: $w_2 = w_1 - 0.1\cdot (-2)/(\sqrt{4}+10^{-8}) = 0 - 0.1\cdot (-2)/2 = 0.1$
- $f(w_2) = (0.1-3)^2 = 8.41$

Iteration $t=2$:
- $w_2=0.1$, $g_2=2(0.1-3)=-5.8$, $|g_2|=5.8>C$
- $u_2=-5.8/(5.8/2)=-2$
- $m_2=0.9(-0.2)+0.1(-2)= -0.18-0.2=-0.38$
- $v_2=0.999\cdot 0.004+0.001\cdot 4 = 0.003996+0.004=0.007996$
- $\hat{m}_2 = -0.38/(1-0.9^2) = -0.38/0.19 = -2$, $\hat{v}_2 = 0.007996/(1-0.999^2) = 0.007996/0.001999 = 4$
- $w_3 = 0.1 - 0.1\cdot (-2)/2 = 0.2$
- $f(w_3)=(0.2-3)^2=7.84$

Iteration $t=3$:
- $w_3=0.2$, $g_3 = 2(0.2-3)=-5.6$, $|g_3|>C$
- $u_3=-2$
- $m_3=0.9(-0.38)+0.1(-2)=-0.342-0.2=-0.542$
- $v_3=0.999\cdot 0.007996+0.001\cdot 4=0.007988004+0.004=0.011988004$
- $\hat{m}_3=-0.542/(1-0.9^3)=-0.542/0.271=-2$, $\hat{v}_3=0.011988004/(1-0.999^3)=0.011988004/0.002997001=4$
- $w_4 = 0.2 - 0.1\cdot (-2)/2 = 0.3$
- $f(w_4)=(0.3-3)^2=7.29$

\section*{Assumptions}
We analyze the clipped method under standard smooth non-convex assumptions.

\begin{assumption}[Smoothness and boundedness]
\label{ass:smooth}
- $f:\R^d\to\R$ is differentiable and $L$-smooth: for all $x,y\in\R^d$,
$\|\nabla f(x)-\nabla f(y)\|_2 \le L \|x-y\|_2$.
- $f$ is bounded below: $\inf_{x\in\R^d} f(x) \ge f_\inf > -\infty$.
\end{assumption}

\begin{assumption}[Algorithmic parameters]
\label{ass:params}
- Clipping threshold $C>0$, stepsize $\eta>0$.
- For the main theorem, we analyze the special case $\beta_1=\beta_2=0$ (Clipped Gradient Descent, CGD). This is a special case of the proposed algorithm (set $\beta_1=\beta_2=0$ and remove bias-corrections).
\end{assumption}

\section*{Main Convergence Theorem}
\begin{theorem}[Global convergence of Clipped Gradient Descent]
\label{thm:main}
Let Assumptions \ref{ass:smooth}--\ref{ass:params} hold and consider CGD
$x_{t+1} = x_t - \eta \, u_t$ with $u_t=\operatorname{clip}_C(\nabla f(x_t))$ and stepsize $\eta \le 1/L$.
Define the sets of iterations
$S_{\le C} \triangleq \{t\in\{1,\dots,T\}:\ \|\nabla f(x_t)\|_2 \le C\}$ and $S_{>C} \triangleq \{1,\dots,T\}\setminus S_{\le C}$.
Then for all $T\ge 1$:
\begin{align}
\text{(i) Monotonic descent: } \quad & f(x_{t+1}) \le f(x_t) \quad \forall t\ge 1. \label{eq:thm-descent}\\
\text{(ii) Large-gradient steps are finite: } \quad & |S_{>C}| \le \frac{2\,(f(x_1)-f_\inf)}{\eta\, C^2}. \label{eq:thm-large}\\
\text{(iii) Summability on small-gradient steps: } \quad &
\sum_{t\in S_{\le C}} \|\nabla f(x_t)\|_2^2 \le \frac{2\,(f(x_1)-f_\inf)}{\eta}. \label{eq:thm-small-sum}\\
\text{(iv) Stationarity rate: } \quad &
\min_{1\le t\le T} \|\nabla f(x_t)\|_2^2 \le \frac{2\,(f(x_1)-f_\inf)}{\eta\,\max\{1,\ T - \frac{2\,(f(x_1)-f_\inf)}{\eta\, C^2}\}}. \label{eq:thm-rate}
\end{align}
In particular, for all $T > \frac{2\,(f(x_1)-f_\inf)}{\eta\, C^2}$,
$\min_{1\le t\le T} \|\nabla f(x_t)\|_2^2 \le \frac{2\,(f(x_1)-f_\inf)}{\eta\, (T - \frac{2\,(f(x_1)-f_\inf)}{\eta\, C^2})} = \mathcal{O}(1/T)$.
\end{theorem}

\section*{Theoretical Properties}
- Stability via clipping: The step direction is always bounded in norm by $C$, ensuring $\|x_{t+1}-x_t\|_2 \le \eta C$. This controls the local Lipschitz error term.
- Hyperparameter sensitivity:
  - If $\eta \le 1/L$, monotone descent holds.
  - Larger $C$ leads to fewer ``large-gradient'' steps but potentially larger oscillations; smaller $C$ yields stronger regularization of steps.
- Comparison to un-clipped GD: Setting $C=\infty$ recovers classical GD and the descent lemma proof gives the standard non-convex rate $\min_{t\le T}\|\nabla f(x_t)\|_2^2 = \mathcal{O}(1/T)$ for $\eta\le 1/L$. Clipping preserves this rate and additionally bounds the number of large-gradient excursions.

\section*{Discussion}
- Relation to Clipped Adam: The proposed practical algorithm uses momentum and second-moment adaptivity with clipping applied to raw gradients. The above theorem establishes a baseline non-asymptotic guarantee for the special case without momentum/adaptivity. This covers the core effect of clipping and smoothness. Extensions to momentum and Adam-style preconditioning typically introduce multiplicative constants depending on $\beta_1,\beta_2$ and require potential-function arguments; these do not change the asymptotic $\mathcal{O}(1/T)$ rate for the minimal gradient-norm squared under standard small-stepsize conditions.
- Deterministic vs. stochastic: The above analysis is deterministic. In stochastic settings, expectations can be applied throughout. With clipping, one also controls heavy-tailed noise; the expectation effectively replaces each occurrence of $f(x_t)$ and $\nabla f(x_t)$ by their expectations, and the same inequalities hold in expectation because the proof is linear in these quantities.

\section*{Preliminaries (Definitions and Lemmas)}
\begin{definition}[Clipping factor]
For $g\in\R^d$ and $C>0$, define the scalar clipping factor
\[
s(g) \;\triangleq\; \min\left\{1,\ \frac{C}{\|g\|_2}\right\}\in (0,1] \,,
\]
and the clipped gradient $u = \operatorname{clip}_C(g) = s(g)\, g$.
\end{definition}

\begin{lemma}[Smoothness descent lemma]
\label{lem:smooth}
Under Assumption~\ref{ass:smooth}, for any $x,y\in\R^d$,
\[
f(y) \;\le\; f(x) + \langle \nabla f(x), y-x\rangle + \frac{L}{2}\,\|y-x\|_2^2.
\]
Proof:
\begin{align*}
\text{[Step 1]}&\ \text{Define }\phi(\tau)=f\big(x+\tau(y-x)\big),\ \tau\in[0,1].\ \text{Then } \phi'(0)=\langle \nabla f(x), y-x\rangle.\\
\text{[Step 2]}&\ \text{By the fundamental theorem of calculus, } f(y)-f(x)=\int_0^1 \phi'(\tau)\,d\tau.\\
\text{[Step 3]}&\ \phi'(\tau) = \langle \nabla f\big(x+\tau(y-x)\big), y-x\rangle.\\
\text{[Step 4]}&\ \phi'(\tau) = \langle \nabla f(x), y-x\rangle + \langle \nabla f\big(x+\tau(y-x)\big)-\nabla f(x), y-x\rangle.\\
\text{[Step 5]}&\ \text{Take absolute value of the second term and use Cauchy-Schwarz: }\\
&\ |\langle \nabla f(x+\tau(y-x))-\nabla f(x), y-x\rangle|
\le \|\nabla f(x+\tau(y-x))-\nabla f(x)\|_2 \cdot \|y-x\|_2.\\
\text{[Step 6]}&\ \text{Apply $L$-smoothness: } \|\nabla f(x+\tau(y-x))-\nabla f(x)\|_2 \le L\,\tau\,\|y-x\|_2.\\
\text{[Step 7]}&\ \Rightarrow \phi'(\tau) \le \langle \nabla f(x), y-x\rangle + L\,\tau\,\|y-x\|_2^2.\\
\text{[Step 8]}&\ \text{Integrate over $\tau\in[0,1]$: } f(y)-f(x) \le \langle \nabla f(x), y-x\rangle + \int_0^1 L\,\tau\,\|y-x\|_2^2\, d\tau.\\
\text{[Step 9]}&\ \int_0^1 L\,\tau\, d\tau = L/2 \ \Rightarrow\ f(y)\le f(x) + \langle \nabla f(x), y-x\rangle + \frac{L}{2}\|y-x\|_2^2.\qedhere
\end{align*}
\end{lemma}

\begin{lemma}[Properties of norm clipping]
\label{lem:clip}
For $g\in\R^d$, $u=\operatorname{clip}_C(g)=s(g)g$ with $s(g)=\min\{1, C/\|g\|_2\}$:
\begin{align}
\text{[Step 10]}&\ \ \ \ \|u\|_2 \le C. \label{eq:clip-norm}\\
\text{[Step 11]}&\ \ \ \ \langle g, u \rangle = s(g)\,\|g\|_2^2 = \min\{\|g\|_2, C\}\,\|g\|_2. \label{eq:clip-inner}\\
\text{[Step 12]}&\ \ \ \ \|u\|_2^2 = s(g)^2\,\|g\|_2^2 = \min\{\|g\|_2^2, C^2\}. \label{eq:clip-normsq}
\end{align}
Proof: Immediate from $u=s(g)g$ and the definition of $s(g)$. For \eqref{eq:clip-norm}, $\|u\|_2=s(g)\|g\|_2\le C$ by the choice of $s(g)$. For \eqref{eq:clip-inner}, $\langle g,u\rangle = s(g)\langle g,g\rangle=s(g)\|g\|_2^2$. For \eqref{eq:clip-normsq}, $\|u\|_2^2=s(g)^2\|g\|_2^2$. \qed
\end{lemma}

\section*{Main Proof (Step-by-Step Derivation)}
We prove Theorem~\ref{thm:main}.

\noindent
Proof:
Let $g_t=\nabla f(x_t)$ and $u_t=\operatorname{clip}_C(g_t)$. Consider the CGD update $x_{t+1}=x_t-\eta u_t$ with $\eta\le 1/L$.

\noindent
\textbf{One-step descent.}
Apply Lemma~\ref{lem:smooth} with $x=x_t$ and $y=x_{t+1}=x_t-\eta u_t$:
\begin{align}
\text{[Step 13]}\ \ f(x_{t+1})
&\le f(x_t) + \langle \nabla f(x_t), -\eta u_t\rangle + \frac{L}{2}\,\|\eta u_t\|_2^2. \nonumber\\
\text{[Step 14]}\ \ 
&= f(x_t) - \eta \langle g_t, u_t\rangle + \frac{L\eta^2}{2}\,\|u_t\|_2^2. \label{eq:one-step}
\end{align}
By Lemma~\ref{lem:clip}, \eqref{eq:clip-inner} and \eqref{eq:clip-normsq}:
\begin{align}
\text{[Step 15]}\ \ \langle g_t, u_t\rangle &= \min\{\|g_t\|_2, C\}\,\|g_t\|_2, \\
\text{[Step 16]}\ \ \|u_t\|_2^2 &= \min\{\|g_t\|_2^2, C^2\}.
\end{align}
Plugging these identities into \eqref{eq:one-step} yields
\begin{align}
\text{[Step 17]}\ \ f(x_{t+1})
&\le f(x_t) - \eta \,\min\{\|g_t\|_2, C\}\,\|g_t\|_2 + \frac{L\eta^2}{2}\,\min\{\|g_t\|_2^2, C^2\}. \label{eq:master}
\end{align}

\noindent
\textbf{Case analysis.}
Define $A_t\triangleq \{\|g_t\|_2 \le C\}$ and $B_t\triangleq \{\|g_t\|_2 > C\}$.
- If $A_t$ holds, then $\min\{\|g_t\|_2, C\}=\|g_t\|_2$ and $\min\{\|g_t\|_2^2, C^2\}=\|g_t\|_2^2$, so
\begin{align}
\text{[Step 18]}\ \ f(x_{t+1})
&\le f(x_t) - \eta\,\|g_t\|_2^2 + \frac{L\eta^2}{2}\,\|g_t\|_2^2 \nonumber\\
\text{[Step 19]}\ \ 
&= f(x_t) - \eta\left(1-\frac{L\eta}{2}\right)\|g_t\|_2^2 \nonumber\\
\text{[Step 20]}\ \ 
&\le f(x_t) - \frac{\eta}{2}\,\|g_t\|_2^2 \qquad \text{since }\eta\le \frac{1}{L}\Rightarrow 1-\frac{L\eta}{2}\ge \frac{1}{2}. \label{eq:small-grad}
\end{align}
- If $B_t$ holds, then $\min\{\|g_t\|_2, C\}=C$ and $\min\{\|g_t\|_2^2, C^2\}=C^2$, so
\begin{align}
\text{[Step 21]}\ \ f(x_{t+1})
&\le f(x_t) - \eta\, C\,\|g_t\|_2 + \frac{L\eta^2}{2}\,C^2 \nonumber\\
\text{[Step 22]}\ \ 
&\le f(x_t) - \eta\, C^2 + \frac{L\eta^2}{2}\,C^2 \qquad \text{since $\|g_t\|_2>C\Rightarrow C\,\|g_t\|_2 \ge C^2$} \nonumber\\
\text{[Step 23]}\ \ 
&= f(x_t) - \eta C^2\left(1-\frac{L\eta}{2}\right) \nonumber\\
\text{[Step 24]}\ \ 
&\le f(x_t) - \frac{\eta C^2}{2} \qquad \text{since }\eta\le \frac{1}{L}. \label{eq:large-grad}
\end{align}

\noindent
\textbf{Monotonicity.}
Combining \eqref{eq:small-grad} and \eqref{eq:large-grad} shows $f(x_{t+1})\le f(x_t)$ for all $t$ [Step 25], proving \eqref{eq:thm-descent}.

\noindent
\textbf{Bounding the number of large-gradient steps.}
Summing \eqref{eq:large-grad} over $t\in S_{>C}$,
\begin{align}
\text{[Step 26]}\ \ f(x_{T+1})
&\le f(x_1) - \sum_{t\in S_{>C}} \frac{\eta C^2}{2} \nonumber\\
\text{[Step 27]}\ \ 
&\le f_\inf \qquad \text{(since $f$ is bounded below by $f_\inf$).} \nonumber
\end{align}
Rearranging,
\begin{align}
\text{[Step 28]}\ \ \frac{\eta C^2}{2}\,|S_{>C}| \le f(x_1)-f_\inf
\quad\Rightarrow\quad
|S_{>C}| \le \frac{2\,(f(x_1)-f_\inf)}{\eta\, C^2}, \label{eq:S-large}
\end{align}
which is \eqref{eq:thm-large}.

\noindent
\textbf{Summability of small-gradient steps.}
Summing \eqref{eq:small-grad} over $t\in S_{\le C}$ gives
\begin{align}
\text{[Step 29]}\ \ f(x_{T+1})
&\le f(x_1) - \sum_{t\in S_{\le C}} \frac{\eta}{2}\,\|g_t\|_2^2 \nonumber\\
\text{[Step 30]}\ \ 
&\le f_\inf \nonumber
\end{align}
and thus
\begin{align}
\text{[Step 31]}\ \ \sum_{t\in S_{\le C}} \|g_t\|_2^2 \le \frac{2\,(f(x_1)-f_\inf)}{\eta}, \label{eq:S-small}
\end{align}
which is \eqref{eq:thm-small-sum}.

\noindent
\textbf{Stationarity rate.}
If $T \le |S_{>C}|$, the bound \eqref{eq:thm-large} already limits how long the large-gradient regime can persist [Step 32]. If $T > |S_{>C}|$, then $|S_{\le C}|\ge T - |S_{>C}|$ and
\begin{align}
\text{[Step 33]}\ \ \min_{1\le t\le T} \|g_t\|_2^2
&\le \frac{1}{|S_{\le C}|} \sum_{t\in S_{\le C}} \|g_t\|_2^2 \nonumber\\
\text{[Step 34]}\ \ 
&\le \frac{2\,(f(x_1)-f_\inf)}{\eta\, |S_{\le C}|}
\le \frac{2\,(f(x_1)-f_\inf)}{\eta\, \left(T - \frac{2\,(f(x_1)-f_\inf)}{\eta\, C^2}\right)}, \label{eq:rate-final}
\end{align}
using \eqref{eq:S-small} and \eqref{eq:S-large} in [Step 34]. This proves \eqref{eq:thm-rate}. \qed

\section*{Rate Derivation (With Constants)}
From \eqref{eq:rate-final}, for any $T > \frac{2\,(f(x_1)-f_\inf)}{\eta\, C^2}$,
\[
\min_{1\le t\le T} \|\nabla f(x_t)\|_2^2 \;\le\; \frac{2\,(f(x_1)-f_\inf)}{\eta\, \left(T - \frac{2\,(f(x_1)-f_\inf)}{\eta\, C^2}\right)}
\;\le\; \frac{4\,(f(x_1)-f_\inf)}{\eta\, T},
\]
where the last inequality uses $T - a \ge T/2$ provided $T\ge 2a$ [Step 35], with $a=\frac{2\,(f(x_1)-f_\inf)}{\eta\, C^2}$. Thus, an explicit constant-rate bound is
\[
\min_{1\le t\le T} \|\nabla f(x_t)\|_2^2 \;\le\; \frac{4\,(f(x_1)-f_\inf)}{\eta}\cdot \frac{1}{T}
\quad \text{for all}\quad T \ge \frac{4\,(f(x_1)-f_\inf)}{\eta\, C^2}.
\]

\section*{Special Cases}
- No clipping ($C\to\infty$): Then $u_t=\nabla f(x_t)$, and \eqref{eq:master} with the $A_t$-case alone yields
$f(x_{t+1}) \le f(x_t) - \eta(1-\frac{L\eta}{2})\|\nabla f(x_t)\|_2^2$. For $\eta\le 1/L$,
$\sum_{t=1}^T \|\nabla f(x_t)\|_2^2 \le \frac{2\,(f(x_1)-f_\inf)}{\eta}$, implying
$\min_{1\le t\le T}\|\nabla f(x_t)\|_2^2 \le \frac{2\,(f(x_1)-f_\inf)}{\eta\, T}$.
- With momentum and adaptivity (Clipped Adam): The practical algorithm includes $m_t,v_t$ and bias-corrections. Although the full proof requires additional potential-function arguments and coordinate-wise bookkeeping, the key descent mechanism remains governed by the clipped direction $u_t$. Under sufficiently small effective stepsizes (componentwise), analogous inequalities to \eqref{eq:small-grad}--\eqref{eq:large-grad} hold with constants scaled by factors depending on $(1-\beta_1)^{-1}$ and the preconditioner bounds induced by $\epsilon$ and the second moment. The asymptotic $\mathcal{O}(1/T)$ rate for $\min_{t\le T}\|\nabla f(x_t)\|_2^2$ is preserved up to multiplicative constants.

\end{document}